Concentrated equity portfolios (sector and thematic funds, high-conviction books, systematic baskets on narrow universes) have a statistical problem that diversified quantitative portfolios do not have. With few names and one dominant common factor, stock-selection alpha is hard to measure, easy to confuse with beta and momentum, and easy for a risk overlay to manufacture. Using point-in-time sell-side target prices, I show that a first-generation regime-timed strategy owed its development Calmar ratio of 1.77 to market exposure. Selection added little. The replacement, First-Mover Flow (FMF), is a zero-beta book built on \emph{leader minus herd}, a signal that separates revisions breaking an information silence from copycat revisions after earnings clusters. Traded in overlapping tranches matched to the signal horizon, FMF has a net Sharpe ratio of 1.13 over 2017--2021 and a factor-adjusted alpha of 7.4\% a year ($t = 2.98$). On 30 stocks over a never-seen period (2024--2026), the frozen, pre-registered specification has a Sharpe ratio of 1.05 and a Calmar ratio of 1.32 at zero beta, and leader minus herd replicates. I then test nine objections. The edge is timing, not a hindsight tilt (causal timing book $t = 2.41$; $t = 2.21$ out of sample). It survives execution delay, borrow fees and costs (breakeven 234\,bp). 99\% of about 300 alternative specifications have a positive Sharpe ratio. A Bayesian random-effects meta-analysis over three independent samples gives a 0.96 probability that the signal works in the next universe. A hybrid GARCH-diffusion generator passes the discriminative and memorisation tests that TimeGAN fails; with it, the single in-sample Calmar ratio becomes a posterior-predictive distribution, and momentum crashes appear as the binding stress.
